% Results and discussion: what the measurements mean for the design (issue #63).
% Cycle counts are from results/phase2/ (#69, #64), resources and clock from
% results/phase2/synth/synth_summary.csv and the Yosys logs next to it (#75).

\subsection{Discussion}
\textbf{The design meets its contract.} Every one of the 106 Phase~1 runs, 112 $K \neq N$ runs and 26 tiled products took exactly the number of cycles that the specification states, on both designs, and the cycle model of Section~III-E reproduces every measured count. The model can therefore be trusted for sizes we did not simulate, which is how Fig.~\ref{fig:cycles} extends the measurements to $64 \times 64 \times 64$.

\textbf{Where the speed-up comes from.} The array is faster because 16 multipliers work at once, not because it runs at a higher clock or uses more logic. Its busy-cycle advantage grows from 5.3 at $k_{\mathrm{len}} = 4$ towards the peak of $N^2 = 16$ as the run gets longer, because the $2N$ cycles of clear, fill, drain and capture are paid once per run. Both designs have the same kind of critical path, operand selection feeding a combinational multiply-accumulate, so the array clocks 9\,\% faster rather than slower, and its time advantage, 5.8 to 11.6 for single runs, is slightly larger than its cycle advantage. In area the array is not the larger design. The flip-flop counts follow the storage each design needs: both hold the same 1024 bits of operands and a 512-bit result store, which in the array is the output buffer, and the array's extra 670 flip-flops are mostly its PE accumulators and forwarding registers. The baseline uses 49\,\% more LUTs, mostly because it picks one operand pair and one result element by run-time index in every cycle; Yosys maps it to 362 \texttt{PFUMX} and 45 \texttt{L6MUX21} wide-multiplexer cells against 219 and none for the array. Per LUT and per flip-flop the array therefore does 17 and 8.2 times as much work, and only per multiplier is the baseline ahead. Within the device's 28 multipliers, $N = 4$ is the largest power-of-two array that fits, and Section~V-C showed why $N$ should be a power of two.

\textbf{The host interface is the bottleneck.} Across the 26 tiled products the array is busy for only 22.9 to 30.0\,\% of the host procedure, so the end-to-end speed-up, 1.2 to 3.5, is far below the busy-cycle one. The same cycle accounting shows what a wider interface would buy. With ports that write one row of $A$ and one column of $B$ per cycle and read one row of $C$ per cycle, a tile would load in $K$ cycles instead of $NK$ and be read in $N$ cycles instead of $N^2$. The $8 \times 8 \times 8$ product would then take 116 cycles instead of 260 on the array and 564 instead of 708 on the baseline, raising the end-to-end speed-up from 2.7 to 4.9. Double-buffered operand storage, which loads the next tile while the current one computes, would hide most of the remaining loading time, as the TPU does with its weight tiles~\cite{jouppi2017tpu}.

\textbf{Fill, drain and padding.} Inside the array, the wavefront limits a PE to $k_{\mathrm{len}} / (k_{\mathrm{len}} + 2N - 2)$ of the compute steps, at most 72.7\,\% with $K_{\max} = 16$. A larger $K_{\max}$ would raise this ceiling, to 91.4\,\% at 64, at the cost of four times the operand storage in each skew buffer. Zero padding is the other loss: the edge tiles of a $5 \times 7 \times 6$ product cost as much time as full ones, so its utilization is 21.9\,\% against 50.0\,\% for $8 \times 8 \times 8$.
